\documentclass[11pt]{article}
\usepackage[a4paper,margin=1in]{geometry}
\usepackage[T1]{fontenc}
\usepackage{lmodern}
\usepackage{amsmath,amssymb,amsthm,mathtools}
\usepackage{microtype}
\usepackage[numbers]{natbib}
\usepackage{hyperref}
\newtheorem{theorem}{Theorem}
\newtheorem{lemma}[theorem]{Lemma}
\newtheorem{proposition}[theorem]{Proposition}
\newtheorem{corollary}[theorem]{Corollary}
\newcommand{\bond}{\operatorname{bond}}
\newcommand{\poly}{\operatorname{poly}}
\title{Packing circuits below the corank: a kernel for binary matroids and edge-disjoint bond packing}
\author{Rong Zhou\\ \small Independent researcher, ORCID \href{https://orcid.org/0009-0004-8825-6103}{0009-0004-8825-6103}}
\date{}
\begin{document}
\maketitle

\begin{abstract}
For a matroid $M$, let $\nu(M)$ be the maximum number of pairwise
element-disjoint circuits and let $\delta(M)=r^*(M)-\nu(M)$, where $r^*$ is
the corank.  We give a polynomial-time kernelization for deciding
$\delta(M)\leq k$ on binary matroids supplied by a binary representation.
Its element bound is $2^{O(4^k)}$; this is a double-exponential bound.
The argument combines a counting lemma valid for every matroid, an
exchange rule for repeated binary circuit profiles, and contraction of
series pairs.  For simple graphs we also give a direct treatment of
edge-disjoint bonds: every two-connected block $H$ on at least three
vertices satisfies $\tau(H)\leq4\delta_{\bond}(H)$.  A false-twin
surplus rule yields a $2^{O(k)}$-vertex kernel using independent block
copies.
\end{abstract}

\section{Introduction}
Every circuit of $M$ contains an element outside any fixed basis, so
$\nu(M)\leq r^*(M)$.  We study the amount by which the maximum packing
falls below this elementary upper bound.  For the graphic case, where
$r^*$ is the cyclomatic number and circuits are cycles,
Harant et al.~\citep[Theorem~3]{Harant2010} proved that for each fixed
$k$ the relevant reduced graphs form a finite family, which yields
fixed-parameter tractability for this graphic case.  The present binary
representation permits circuits beyond graphic matroids.
Rautenbach and Regen~\citep{RautenbachRegen2012} study the distinct
vertex-disjoint cycle variant.  The dual of a graphic matroid gives a
second motivating case: its circuits are graph bonds, and its corank is
$|V|-c(G)$.  Every nonempty cut contains a bond, so replacing each
member of an edge-disjoint cut packing by a contained bond preserves
the packing size; the two maximum packing numbers are equal.  Cut
packing was studied by Caprara, Panconesi, and Rizzi~\citep{Caprara2004}.

Put
\begin{equation}\label{eq:g}
 g(k)=2k+2^{2k}+2k\,2^{2^{2k}}2^{2k}\qquad(k\geq0).
\end{equation}
In particular $g(k)=2^{4^k+O(k)}$ for positive $k$.

\begin{theorem}[Binary-matroid kernel]\label{thm:main}
Given a binary matrix $A$ representing $M$ and
$k\in\mathbb Z_{\geq0}$, there is a polynomial-time algorithm, measured
in the input encoding length $L$, that outputs an equivalent instance
$(A',k')$ with $k'\leq k$ and at most $\max\{3,g(k)\}$ elements.
After row reduction, $A'$ has $O(\max\{3,g(k)\}^2)$ bits.  The problem
is fixed-parameter tractable.  For $k<0$, the algorithm outputs a
fixed NO instance.
\end{theorem}

\begin{theorem}[Bond packing]\label{thm:bond-main}
Let $G$ be a finite simple graph, and let $\nu_{\bond}(G)$ be the maximum
number of pairwise edge-disjoint bonds.  With
$\delta_{\bond}(G)=|V(G)|-c(G)-\nu_{\bond}(G)$, deciding
$\delta_{\bond}(G)\leq k$ is fixed-parameter tractable and has a
$2^{O(k)}$-vertex kernel.  Every two-connected block $H$ on at least
three vertices satisfies $\tau(H)\leq4\delta_{\bond}(H)$.
\end{theorem}

The binary-matroid proof does not rely on recognizing an optimal packing:
it finds a small set of exceptional nonbasis elements by vertex cover.
Since a cographic matroid is binary, the FPT assertion of
Theorem~\ref{thm:bond-main} also follows from Theorem~\ref{thm:main}
via a polynomial-time binary representation.  The direct graph proof
gives the vertex-cover bound and the smaller vertex kernel: in a YES
instance at most $k$ nontrivial blocks survive, each with at most
$4k\,2^{4k-1}$ vertices.

\section{Preliminaries}\label{sec:prelim}
All matroids are finite.  A circuit is a minimal dependent set, and
$B$ denotes a basis.  Write $B^*=E(M)\setminus B$, so
$|B^*|=r^*(M)=|E(M)|-r(M)$.  For $e\in B^*$ let $C(e,B)$ denote its
fundamental circuit.  Loops and coloops are allowed.  The binary input
representation permits Gaussian elimination, construction of a basis,
contraction, and computation of fundamental circuits in polynomial time.

A \emph{cycle} of a binary matroid means a vector of its binary circuit
space, identified with its support.  This includes the empty cycle.
Addition is symmetric difference.  The fundamental circuits relative
to $B$ form a vector-space basis: each contains exactly one element of
$B^*$, namely its index.  Every nonempty cycle is an element-disjoint
union of circuits.  Contraction by a set projects the cycle space onto
the coordinates outside that set.  These facts follow as well by writing
a binary representation in basis form and taking its nullspace.

We use two elementary reductions.  (R1) Delete a coloop.  (R2) If
$\{x,y\}$ is a two-element cocircuit, contract one member.  Both preserve
$\delta$.  For R1, no circuit uses the coloop and both $r^*$ and $\nu$
are unchanged.  For R2, binary circuit--cocircuit orthogonality says
that every cycle contains $x$ if and only if it contains $y$.
Projection deleting the $x$ coordinate is injective on the cycle space:
an element of its kernel would be the singleton cycle $\{x\}$, which
cannot contain $x$ without $y$.  Membership of $x$ is recovered from
membership of $y$, so support containment and disjointness are preserved
in both directions.  The projection therefore bijects circuits and
their disjoint packings.  It preserves the circuit-space dimension
$r^*$ and the optimum $\nu$.  The pair may lie inside a larger
component; R2 contracts one member.

For a graph, a bond is a minimal nonempty edge cut.  Equivalently, in a
connected component it is a cut $\partial(S)$ for which both shores
induce connected graphs.  The two definitions agree: if one shore is
disconnected, the edges leaving one of its components form a proper
nonempty cut subset; conversely, if both shores are connected, omitting
any cut edge reconnects the shores, so no proper subset is a cut.
We write $\tau(G)$ for the minimum vertex-cover size.  A two-connected
block below always has at least three vertices; bridge blocks are treated
separately.

\section{A counting lemma for arbitrary matroids}
\begin{lemma}\label{lem:count}
For any matroid $M$, any basis $B$, and any maximum circuit packing
$\mathcal P$, let $q$ be the number of packed circuits containing at
least two members of $B^*$.  Then $q\leq\delta(M)$.  The other packed
circuits are pairwise disjoint fundamental circuits $D_i=C(e_i,B)$.
If $F=B^*\setminus\{e_i\}$, then
\begin{equation}\label{eq:count}
 |F|=\delta(M)+q\leq2\delta(M).
\end{equation}
\end{lemma}
\begin{proof}
A circuit cannot lie in the independent set $B$, so every packed
circuit contains at least one member of $B^*$.  Put
$a_C=|C\cap B^*|$.  Since the packing is element-disjoint,
$\sum_{C\in\mathcal P}a_C\leq|B^*|=\nu(M)+\delta(M)$.
Since every $a_C\geq1$ and $q$ of them are at least two, the same sum
is at least $\nu(M)+q$.  Thus $q\leq\delta(M)$.  If $a_C=1$ with
$C\cap B^*=\{e\}$, then $C\subseteq B\cup\{e\}$ and uniqueness of
the circuit in $B\cup\{e\}$ gives $C=C(e,B)$.  There are
$\nu(M)-q$ such circuits, with distinct indexing elements and disjoint
supports.  Subtracting their indexing elements from $B^*$ gives
$|F|=r^*(M)-(\nu(M)-q)=\delta(M)+q$.
\end{proof}

\section{An exchange lemma for repeated binary profiles}\label{sec:exchange}
Fix a basis $B$ of a binary matroid and \emph{any} pairwise disjoint
family $D_i=C(e_i,B)$ indexed by $I$; this family need not be optimal.
Write
\[
 D_i=\{e_i\}\mathbin{\dot\cup}P_i,\quad
 F=B^*\setminus\{e_i:i\in I\},\quad
 P_0=B\setminus\bigcup_{i\in I}P_i,
\]
and, for $f\in F$,
\[
 C(f,B)=\{f\}\mathbin{\dot\cup}R_f
             \mathbin{\dot\cup}\bigsqcup_{i\in I}A_{fi},
 \qquad R_f\subseteq P_0,\quad A_{fi}\subseteq P_i.
\]
For $x\in P_i$ define its full signature
$s_i(x)=\{f\in F:x\in A_{fi}\}$.  The \emph{profile} of $D_i$ is the
multiset of these signatures for $x\in P_i$, with the labels of $F$
fixed.  Keeping the entire joint signature, rather than only the
individual cardinalities $|A_{fi}|$, is essential.

\begin{lemma}[Exchange or H1]\label{lem:exchange}
If $T\subseteq I$ consists of $t$ circuits with the same profile and
$t>|F|$, then for any $i_0\in T$,
\[
 \nu(M/D_{i_0})=\nu(M)-1,
 \qquad \delta(M/D_{i_0})=\delta(M).
\]
\end{lemma}
\begin{proof}
Put $h=|F|$ and $D=D_{i_0}$.  First assume $h\geq1$, so after the
contraction there are $s=t-1\geq h\geq1$ retained circuits of type
$T$.  Equality of signature multisets supplies a common position set
$L$ and bijections $\phi_i:L\to P_i$ such that the signature at a
given position is the same in every $P_i$.  For $Z\subseteq F$ put
$a_i(Z)=\mathop{\triangle}_{f\in Z}A_{fi}$.  Under $\phi_i$, all sets
$a_i(Z)$ correspond to the same subset $a(Z)\subseteq L$; in
particular they are simultaneously empty or nonempty.

Every cycle $X$ of $M$ has a unique expression
\begin{equation}\label{eq:cycleexp}
 X=\mathop{\triangle}_{f\in Z}C(f,B)\ \triangle\
   \mathop{\triangle}_{i\in I:u_i=1}D_i,
 \qquad Z=X\cap F,
 \quad u_i=\mathbf 1_{e_i\in X}.
\end{equation}
Only the indicated fundamental circuits meet $D_i$, hence
\begin{equation}\label{eq:local}
 X\cap D_i=\begin{cases}a_i(Z),&u_i=0,\\
                         D_i\setminus a_i(Z),&u_i=1.
          \end{cases}
\end{equation}
When $a(Z)\ne\varnothing$, both alternatives in
\eqref{eq:local} are nonempty: $a_i(Z)\subseteq P_i$ while the other
alternative contains $e_i$.  Also $Z\ne\varnothing$.  A circuit of
this kind meets every $T$-block and uses at least one member of $F$;
call it \emph{crossing}.  A disjoint packing has at most $h$ crossing
circuits, because their nonempty intersections with $F$ are disjoint,
and it cannot simultaneously contain a $T$-block circuit.  If
$a(Z)=\varnothing$, the circuit meets each $T$-block in either nothing
or the whole block.  If it contains a whole block, circuit minimality
forces it to equal that block.  Otherwise call it \emph{neutral}.
Thus crossing, block, and neutral circuits exhaust all circuits.

We next justify that the family of neutral circuits is identical in
$M$ and $N=M/D$, as subsets of $E(M)\setminus\bigcup_{i\in T}D_i$.
Since $D$ is a circuit, $B'=B\setminus P_{i_0}$ is a basis of $N$:
it has $r(M)-(|D|-1)$ elements and its image spans the contraction.
The fundamental circuits in $N$ are $C(f,B)\setminus D$ for $f\in F$
and $D_i$ for $i\ne i_0$.  Therefore \eqref{eq:local} still holds
on every retained block.  Take a circuit of $N$ avoiding all retained
blocks.  On any one retained block, absence of $e_j$ gives $u_j=0$,
and then \eqref{eq:local} gives $a_j(Z)=\varnothing$.  Profile equality
gives $a_{i_0}(Z)=\varnothing$.  Lift its fundamental-circuit
expression to $M$, choosing $u_{i_0}=0$; the lift has no coordinates
on $D$, so it has precisely the original support.  It is a cycle of
$M$.  If it contained a smaller nonempty cycle of $M$, that cycle
would avoid $D$ and project to a smaller nonempty cycle of $N$,
contradicting circuit minimality.  Hence it is a circuit of $M$.
Conversely, a neutral circuit of $M$ projects without changing its
support to a nonempty cycle of $N$.  That cycle contains a circuit of
$N$; this circuit avoids retained blocks and so, by the preceding
argument, is also a circuit of $M$.  Minimality in $M$ forces equality.
The arbitrary sets $R_f$ never affect this argument: all coordinates
outside $D$ are retained unchanged.

Let $\rho$ be the maximum packing number of this common neutral
family.  In $M$, a packing without crossing circuits has at most $t$
block circuits and $\rho$ neutral circuits; one with crossing circuits
has no $T$-block circuits and at most $h+\rho<t+\rho$ members.
The $t$ disjoint block circuits together with a maximum neutral
packing attain $t+\rho$, hence $\nu(M)=t+\rho$.  In $N$ the same
crossing, block, and neutral trichotomy applies to the nonempty
retained family $T\setminus\{i_0\}$ of matching profiles.  It gives
an upper bound $\max\{s,h\}+\rho=s+\rho$, attained by
all $s$ retained blocks and a maximum neutral packing.  Thus
$\nu(N)=s+\rho=\nu(M)-1$.  A circuit has rank $|D|-1$, so
\[
 r^*(N)=(|E(M)|-|D|)-\bigl(r(M)-(|D|-1)\bigr)=r^*(M)-1.
\]
Subtracting packing numbers proves the second assertion.

If $h=0$, all fundamental circuits are the disjoint $D_i$; basis
elements outside them are coloops.  Every nonempty cycle is a union
of some $D_i$, and its minimal nonempty supports are exactly the
$D_i$.  Contracting one $D_i$ removes precisely that circuit and
reduces both $r^*$ and $\nu$ by one.  This includes $t=1$ and a
singleton-loop block.
\end{proof}

The strict threshold matters.  Take the binary matroid on
$E=\{0,1,\ldots,7\}$ with basis $B=\{2,4,6,7\}$ and fundamental
circuits $C(0,B)=\{0,2,4\}$, $C(1,B)=\{1,2,6\}$,
$C(3,B)=\{3,4,7\}$, and $C(5,B)=\{5,6,7\}$.  These four circuits
specify its binary cycle space.  The disjoint block circuits
$\{1,2,6\}$ and $\{3,4,7\}$ have the same profile over
$F=\{0,5\}$, so $t=h=2$.  Contracting either one of these two blocks
changes $(r^*,\nu,\delta)$ from $(4,2,2)$ to $(3,2,1)$.
Thus $t=h$ cannot replace $t>h$.

\section{Kernel for binary matroids}\label{sec:kernel}
Fix any basis $B$.  Make a conflict graph whose vertices are $B^*$:
two vertices $e,f$ are adjacent if $C(e,B)$ and $C(f,B)$ intersect.
An independent set is exactly a family of disjoint fundamental
circuits; its complement is a vertex cover.

\begin{lemma}[A small witness can be found]\label{lem:cover}
If $\delta(M)\leq k$, the conflict graph has a vertex cover of size
at most $2k$.  In time $O(4^k\poly(L))$ one can find such a
cover, or correctly reject $\delta(M)\leq k$.
\end{lemma}
\begin{proof}
Apply Lemma~\ref{lem:count} to a maximum packing relative to the
chosen $B$.  Its good fundamental circuits are disjoint and their
indexing vertices form an independent set; the complementary set
$F$ is a cover of size at most $2\delta(M)\leq2k$.  This proves the
existence statement without requiring the algorithm to find the
maximum packing.  For the algorithm, repeatedly take an uncovered
edge $ef$ and branch on putting $e$ or $f$ into the cover.  Each
branch reduces the remaining budget by one; a search to depth $2k$
has at most $2^{2k+1}-1$ nodes, and each node is processed in
polynomial time.  Exhausting all branches proves that no such cover
exists and hence that the instance is a NO instance.  A found cover
is only a witness for subsequent reductions, not proof of YES.
\end{proof}

Given the cover $F$, let $I=B^*\setminus F$.  We apply R1, R2, and
Lemma~\ref{lem:exchange} whenever their conditions hold.  The
exchange rule compares profile multisets and contracts one whole
block from any class occurring more than $|F|$ times.  R1 removes a
zero coordinate column from the cycle-space representation.  R2
contracts one member of a pair of equal nonzero columns.  Equality of
nonzero cycle-space columns is exactly the two-element cocircuit
condition: the corresponding coordinates agree on every cycle,
so their pair is a cocycle, while neither singleton is a cocycle.
When needed, these columns and profiles are obtained by linear algebra
and sorting.  No exact computation of $\nu$ enters the reductions.

The witness can be maintained.  R1 deletes a basis element because a
nonbasis column in the fundamental-circuit representation is a
nonzero unit column.  Every R2 pair has at least one basis element,
since distinct nonbasis elements have distinct unit columns; choose
that basis element for contraction.  After either operation, remove it
from $B$ and from any affected $P_i$ or $P_0$.  The remaining
fundamental circuits are exactly the old ones with the removed
coordinate deleted, remain disjoint, and the labelled set $F$ is
unchanged.  After H1 contracts $D_i=\{e_i\}\mathbin{\dot\cup}P_i$, remove
$P_i$ from $B$ and $i$ from $I$; again $F$ remains unchanged and the
other blocks are fundamental circuits.  Each operation removes at
least one element.  Thus there are at most $|E(M)|$ operations and no
loop caused by witness recomputation.

\begin{lemma}[Terminal size]\label{lem:size}
If $h=|F|\leq2k$ and none of R1, R2, H1 applies to the maintained
witness, then $|E(M)|\leq g(k)$.
\end{lemma}
\begin{proof}
Partition $E(M)=F\mathbin{\dot\cup}P_0\mathbin{\dot\cup}
\bigsqcup_{i\in I}D_i$.
For $x\in P_i$, its coordinate in a cycle with switches
$(Z,(u_j))$ is $u_i+|s_i(x)\cap Z|\pmod2$; the coordinate of $e_i$
is $u_i$.  For $x\in P_0$ it is $|s_0(x)\cap Z|\pmod2$, where
$s_0(x)=\{f\in F:x\in R_f\}$.  These formulas follow by expanding
\eqref{eq:cycleexp}.  A member of $P_0$ with empty signature is a
coloop and would be removed by R1.  Two members with the same
nonempty signature have identical nonzero columns and would trigger
R2.  Consequently $|P_0|\leq2^h-1$.

Within each $D_i$, $e_i$ supplies the empty signature.  A member of
$P_i$ with empty signature agrees with $e_i$ on every cycle and
would trigger R2.  Two members of $P_i$ with the same nonempty
signature would do the same.  Hence $P_i$ contains at most one
representative of each nonempty subset of $F$ and $|D_i|\leq2^h$.
Its profile is now a subset of the $2^h-1$ possible nonempty
signatures, so there are at most $2^{2^h-1}$ profiles.  Since H1 is
inapplicable, each profile occurs at most $h$ times.  Therefore
\[
 |E(M)|\leq h+(2^h-1)+h\,2^{2^h-1}\,2^h
          \leq2k+2^{2k}+2k\,2^{2^{2k}}\,2^{2k}=g(k).
\]
Here $g$ is a convenient upper bound.  Blocks consisting of a single
loop have the empty profile and are included.  If
$h=0$, no block may remain, R1 removes all of $P_0$, and the terminal
matroid is empty; the displayed bound still holds.
\end{proof}

\begin{proof}[Proof of Theorem~\ref{thm:main}]
For the FPT algorithm, find a cover using Lemma~\ref{lem:cover},
reject if none exists, then apply safe reductions to a terminal
instance.  Lemma~\ref{lem:exchange} and the R1/R2 arguments preserve
$\delta$ at every step, and Lemma~\ref{lem:size} bounds the remaining
element count by $g(k)$.  Enumerating the $2^{r^*}$ cycles of the
terminal binary representation, retaining minimal nonzero supports,
and checking all subsets of those supports for pairwise disjointness
computes $\nu$ exactly in a function of $k$ times a polynomial in
the input size.  It decides the original instance, including the
case where a small cover exists but $\delta>k$.

For a \emph{standard polynomial-time kernelization}, one further
wraps this FPT preprocessing.  If $k<0$, return a fixed NO instance
represented by $(U_{1,3},0)$.  For $k\geq0$ and $n=|E(M)|\leq g(k)$,
row-reduce the representation to full row rank and return it with
$k'=k$.  Otherwise run the cover search and reductions.  Since
$n>g(k)\geq4^k$ for $k\geq1$, the single $4^k\poly(L)$ cover search
costs polynomial time in the original input length; there are at most
$n$ reduction steps, each polynomial in $L$.  For $k=0$, $g(0)=1$
and the search has constant branching.  When the current instance
reaches $g(k)$ elements it may be row-reduced and returned with
$k'=k$.  If the search finds no cover, the branch is NO by
Lemma~\ref{lem:cover}; encode it as $(U_{1,3},0)$.  This matroid has
three parallel elements, $r^*=2$, $\nu=1$, and $\delta=1>0$.
Its three elements explain the
$\max\{3,g(k)\}$ bound in the statement.  The comparison with $g(k)$
is carried out with integers saturated at $n+1$, so an enormous value
of $g(k)$ is never explicitly constructed.  In every return branch
the output matrix has full row rank and at most
$O(\max\{3,g(k)\}^2)$ bits, since its rank is at most its number of
columns.  Thus the wrapper is
polynomial-time and gives a kernel; the exhaustive terminal solver
proves FPT.
\end{proof}

\section{Edge-disjoint bond packing}\label{sec:bonds}
We give a graph-specific route to Theorem~\ref{thm:bond-main}.
We use independent copies of blocks: an articulation vertex shared by
two blocks appears once in each block instance.

\begin{lemma}[Block additivity]\label{lem:blocks}
The bonds of a graph are precisely the bonds of its blocks,
regarded as edge sets of the original graph.  In particular, with
bridge blocks included,
\[
 \nu_{\bond}(G)=\sum_H\nu_{\bond}(H),\qquad
 |V(G)|-c(G)=\sum_H(|V(H)|-1),\qquad
 \delta_{\bond}(G)=\sum_H\delta_{\bond}(H).
\]
Every bridge block has deficit zero.
\end{lemma}
\begin{proof}
It suffices to work in a connected component.  The block--cut
incidence graph is a tree.  A simple cycle cannot enter another block
through an articulation vertex and return without repeating that
vertex.  Hence each cycle lies within one block, and the graphic
cycle space is the direct sum of the cycle spaces of the
blocks, on disjoint edge sets.  Taking orthogonal complements over
$\mathbb F_2$, the cut space is likewise the direct sum of the block
cut spaces.  A bond is a nonzero cut-space support minimal under
inclusion: a nonminimal cut-space support contains a smaller nonzero
cut support, and every nonzero cut support contains a minimal one.
Minimal nonzero supports of a direct sum on disjoint coordinate sets
belong to exactly one summand, proving the bond assertion.  Edge
disjoint packings then split and combine blockwise.  The rank identity
is obtained by adding each block to the block--cut tree: its first
articulation vertex has already been counted and its remaining
$|V(H)|-1$ vertices are new.  A bridge block has one edge, itself
a bond, so both its rank contribution and packing contribution are
one.  Summing gives the three formulas.
\end{proof}

\begin{lemma}[Vertex-cover bound]\label{lem:tau}
If $H$ is a simple two-connected graph on at least three vertices,
then $\tau(H)\leq4\delta_{\bond}(H)$.
\end{lemma}
\begin{proof}
Let $d=\delta_{\bond}(H)$, and choose a maximum disjoint bond
packing $\mathcal Q$ of size $|V(H)|-1-d$.  Fix a spanning tree $T$.
Each bond meets $T$, for otherwise all its edges could be removed
without disconnecting $T$.  If $q$ packed bonds meet $T$ in at least
two edges, then the disjointness of the packing gives
\[
 |V(H)|-1\geq\sum_{Q\in\mathcal Q}|Q\cap E(T)|
       \geq |\mathcal Q|+q=|V(H)|-1-d+q;
\]
thus $q\leq d$.  A packed bond meeting $T$ in exactly one edge $e$
is the fundamental cut of $e$: no other tree edge crosses the cut,
so each component of $T-e$ lies in one shore.  Call these
$|\mathcal Q|-q$ tree edges good, and let $F_T$
consist of the remaining tree edges.  Then
$|F_T|=|V(H)|-1-(|\mathcal Q|-q)=d+q\leq2d$.

The tree path of a graph edge can contain at most one good edge.
Indeed, an edge whose path contains a good edge crosses its
fundamental cut and belongs to the corresponding packed bond; two
good edges on the path would place the graph edge in two disjoint
packed bonds.  Contract each component of the forest $(V(H),F_T)$
to a vertex.  The uncontracted good edges form a tree $Q$.  Every
graph edge either remains inside one contracted vertex or joins
adjacent vertices of $Q$, by the preceding path property.

Consider a contracted vertex consisting of a single original vertex
$v$; it is not incident with $F_T$.  If it had degree at least two in
$Q$, deleting $v$ would separate the corresponding branches of
$Q$, and no graph edge could reconnect them, since graph edges join
only equal or adjacent $Q$ vertices.  This contradicts
two-connectivity.  Hence every singleton contracted vertex has
$Q$-degree at most one.  If two singleton contracted vertices were
adjacent by an edge of $H$, they would be adjacent in $Q$.  Since
both have degree at most one, $Q$ would consist of just these two
vertices; as all vertices of $Q$ are singletons, $H$ would have only
two vertices.  This is excluded.  Therefore every edge of $H$ has
an endpoint incident with $F_T$.  The endpoints of $F_T$ form a
vertex cover of size at most $2|F_T|\leq4d$.
\end{proof}

\begin{lemma}[False-twin surplus]\label{lem:twins}
Let $H$ be a simple connected graph, and let $T$ be an independent
set of $t$ vertices with the same open neighborhood $N$, of size
$d\geq1$.  If $t>d$, deleting any one $v\in T$ preserves
$\delta_{\bond}(H)$.
\end{lemma}
\begin{proof}
For $x\in T$, put $D_x=\partial(\{x\})$.  There is another member
of $T$, and any path segment $a-x-b$ with $a,b\in N$ can be rerouted
as $a-y-b$ through such a member $y$.  Hence $H-x$ is connected,
and $D_x$ is a bond.  The $D_x$ are edge-disjoint and each has $d$
edges.

Let $Q=\partial(S)$ be a bond meeting $D_x$ but unequal to it, with
$x\in S$.  If all of $N$ lay in $S$, $Q$ would not meet $D_x$.  If
all of $N$ lay outside $S$, then $x$ would be isolated in $H[S]$;
connectivity of that shore would force $S=\{x\}$ and $Q=D_x$.
Thus each shore contains a vertex of $N$, and $Q$ meets every
$D_y$, $y\in T$.  Call $Q$ crossing.  In an edge-disjoint packing
there are at most $d$ crossing bonds, because each uses a distinct
edge of any fixed star $D_x$.  A crossing bond prevents the presence
of any $D_y$.  Replacing the crossing bonds by all $t>d$ stars
strictly increases the packing.  If there are no crossing bonds,
the packing may be enlarged to contain every star, because every
other packed bond avoids all star edges.  Consequently some maximum
packing contains all $t$ stars.

Let $\rho$ be the largest packing of bonds avoiding every star edge.
A cut avoiding all star edges keeps $T\cup N$ on one shore.  Deleting
$v$ from that shore preserves its connectivity: occurrences of $v$
on a path can be rerouted through a remaining twin, and the other
shore is unchanged.  Conversely, adding $v$ to the shore containing
$T\setminus\{v\}$ and $N$ preserves both shores' connectivity.
The cut edge set itself is unchanged.  Thus these neutral bonds,
including their edge-disjointness relations, correspond in $H$ and
$H-v$, and have the same packing number $\rho$.  We have
$\nu_{\bond}(H)=t+\rho$.

In $H-v$ the twin class has $t-1\geq d$ members.  Its stars are
still bonds: when $t-1\geq2$ the same rerouting proves this, while
the only case $t-1=1$ has $d=1$ and its star is a bridge.  The same
crossing argument bounds their number by $d$.
Replacing crossing bonds by all $t-1$ stars now does not decrease
packing size, even at equality $t-1=d$.  Hence
$\nu_{\bond}(H-v)=t-1+\rho$.  The path rerouting also shows $H-v$
remains connected, so its rank contribution $|V|-1$ decreases by
one.  The deficit is unchanged.
\end{proof}

For isolated false twins, deleting a vertex decreases both $|V|$
and $c(G)$ by one and changes no bond; its deficit is unchanged.

\begin{lemma}[Size of a reduced block]\label{lem:blocksize}
Let $H$ be a two-connected simple block on at least three vertices,
and suppose no false-twin class in this \emph{block instance} has
size larger than its neighborhood.  Then
$|V(H)|\leq\tau(H)2^{\tau(H)-1}$.
\end{lemma}
\begin{proof}
Take a vertex cover $C$ of size $\tau=\tau(H)$, and put
$I=V(H)\setminus C$.  The set $I$ is independent.  A two-connected
graph on at least three vertices has minimum degree at least two,
so every $x\in I$ has a neighborhood $S\subseteq C$ with $|S|\geq2$.
For each such $S$, the vertices in $I$ with neighborhood $S$ are
contained in a false-twin class, which by assumption has at most
$|S|$ members.  These groups partition $I$; hence
\[
 |V(H)|\leq\tau+
  \sum_{S\subseteq C,\,|S|\geq2}|S|
 =\tau+(\tau2^{\tau-1}-\tau)=\tau2^{\tau-1}.
\]
The identity in the middle counts each vertex of $C$ in exactly
$2^{\tau-1}$ subsets and subtracts the $\tau$ singleton subsets.
\end{proof}

\begin{proof}[Proof of Theorem~\ref{thm:bond-main}]
Use Lemma~\ref{lem:blocks} to split $G$ into independent block
instances, with separate copies of every articulation vertex.
Bridge and isolated-vertex instances have zero deficit and may be
discarded.  Within a current block, apply Lemma~\ref{lem:twins};
if a deletion destroys two-connectivity, replace the block by its
new independent block instances.  The deficit is unchanged by the
surplus rule and additive under the replacement.  Every nontrivial
deletion reduces the sum of block rank contributions
$\sum_H(|V(H)|-1)$ by one; splitting preserves that sum.  Thus at
most $|V(G)|-c(G)$ nontrivial deletions occur.  Block decomposition,
twin grouping, and the graph updates take polynomial time per step.

If $k<0$, reject.  If $k=0$, accept exactly when no nontrivial block
remains: each such block has $\tau\geq2$ and
Lemma~\ref{lem:tau} forces positive integral deficit.  For $k\geq1$,
if more than $k$ nontrivial blocks remain, reject by additivity and
the same positive-deficit observation.  Otherwise every surviving
block in a YES instance has deficit at most $k$.  Lemmas
\ref{lem:tau} and \ref{lem:blocksize} therefore bound its size by
$N(k)=4k\,2^{4k-1}$.  If any reduced block exceeds this size, reject.
For a kernel, output the disjoint union of the at most $k$ surviving
nontrivial block instances, with parameter $k$.  Additivity preserves
the deficit, and the union has at most $kN(k)=2^{O(k)}$ vertices.
A rejected branch outputs the fixed NO instance $(K_3,0)$; when no
nontrivial block remains, output the empty graph with parameter $0$.
The reductions and block construction take polynomial time, with
size comparisons saturated at the input vertex count, so this is a
vertex kernel.

For the FPT decision algorithm, on the at most $k$ remaining blocks,
enumerate bipartitions and
retain cuts with connected shores.  A dynamic program whose state is
the set of edges already used either skips a bond or adds it when its
edge set is disjoint from the state.  It finds a maximum edge-disjoint
packing.  Each block has at most $\binom{N(k)}2$ edges and at most
$2^{N(k)-1}$ distinct shores to inspect, so the dynamic program takes
$2^{O(N(k)^2)}\poly(|V(G)|)$ time.  Sum the exact block deficits
and compare with $k$.  All rejection tests are necessary for a YES
instance, and every other step preserves or exactly computes the
deficit; this proves the FPT assertion.
\end{proof}

The independent-copy qualification is material.  In the graph with
edges $02$, $03$, $04$, $12$, $13$, $14$, $25$, $56$, $62$, a $K_{2,3}$ block and a
triangle meet at vertex $2$.  Vertex $2$ belongs to a false-twin
class within the bipartite block, but has additional neighbors in the
whole graph.  Deleting it from the whole graph changes the deficit
from $2$ to $1$, whereas reducing only its copy inside the bipartite
block preserves both block deficits.

\section{Computational verification}\label{sec:comp}
Scripts, summary data, and reproduction commands are available at
\url{https://github.com/zhourong-bot/binary-matroid-circuit-packing}
(archived at Zenodo, \href{https://doi.org/10.5281/zenodo.22952617}{doi:10.5281/zenodo.22952617}).
The binary-matroid checks processed the complete catalogue of $3,766$
isomorphism classes on at most ten elements and its $134,860$ bases,
as well as constructed and random inputs with at most sixteen elements
and selected cographic $K_{a,b}$ instances with up to $32$ elements.
The kernel pipeline was checked on $4,021$ inputs, each at
$k=\delta$ and $k=\delta-1$, for $8,042$ parameter choices.  The
input deficits $0,1,2,3,4,5,6$ occurred respectively $478,1264,1633,
564,54,21,7$ times.  There were $15,716$ local checks: $14,344$
pipeline reductions and $1,372$ nonminimum-cover H1 targets, with no
observed deficit failure.  Among the terminal instances with
$\delta=1,2,3$, the maximum element counts were respectively
$9,13,20$.  The implementation used brute-force cover enumeration on
these small inputs; this does not benchmark the bounded search tree
of Theorem~\ref{thm:main}.

The graph-side enumeration covered all two-connected simple graphs on
at most nine vertices ($201,727$ graphs).  A separate check of the
cover argument processed $70$ two-connected graphs and $9,588$
spanning trees.  These finite checks test the implementations; they
are not proofs of the lemmas or bounds.

\section{Related work}\label{sec:related}
For graphic matroids, Harant et al.~\citep{Harant2010} already give
finite obstruction sets for each deficit $k$ and fixed-parameter
tractability for edge-disjoint cycle packing below the cyclomatic
number.  Rautenbach and Regen~\citep{RautenbachRegen2012} treat the
vertex-disjoint cycle variant.  Theorem~\ref{thm:main} covers
explicitly represented binary matroids.  For bonds, the cographic
matroid $M^*(G)$ is binary and has
$r^*(M^*(G))=r(M(G))=|V(G)|-c(G)$ and
$\nu(M^*(G))=\nu_{\bond}(G)$.  Hence its FPT conclusion follows
from Theorem~\ref{thm:main}.  The graph-specific contribution is the
direct $\tau\leq4\delta_{\bond}$ proof and the $2^{O(k)}$-vertex
kernel with at most $k$ nontrivial blocks of at most
$N(k)=4k\,2^{4k-1}$ vertices each.

Caprara, Panconesi, and Rizzi~\citep{Caprara2004} study packing
cuts.  Since each nonempty cut contains a bond, their cut-packing
setting and edge-disjoint bond packing have the same maximum packing
number.  Gutin and Mnich~\citep{GutinMnich2022} survey parameterization
above and below guaranteed values.  Related matroid packing and
covering problems appear in~\citep{BernathKiraly2014,Oxley1978,Seymour1980}.

\section{Open problems}
\sloppy
The element bound $2^{O(4^k)}$ leaves open whether the binary-matroid
problem has a single-exponential or polynomial kernel.  For
two-connected graph blocks, the proved cover bound is
$\tau\leq4\delta_{\bond}$; enumeration through nine vertices found
no counterexample to $\tau\leq2\delta_{\bond}$, for which we give no
proof.  A direction not treated here is profile exchange and
compression for matroids representable over $GF(q)$ with $q>2$; the binary
symmetric-difference argument above does not establish that case.

\section*{AI-use disclosure}
Multiple AI systems assisted with candidate arguments, proof review,
finite-instance checks, and drafting.  The author takes responsibility
for the manuscript.

\bibliographystyle{plainnat}
\bibliography{refs}
\end{document}
